\section{master服务器操作}
master执行所有命名空间操作。此外，它还管理整个系统中的chunk replica：
负责做出放置决策、创建新的chunk及其replica，并协调各种全系统范围的
活动，以保持chunk具有完整replica、平衡所有chunkserver之间的负载，并回收
未使用的存储空间。下面将分别讨论这些主题。

\subsection{命名空间管理与锁}
许多master操作可能耗时很长。例如，\textbf{snapshot}操作必须撤销\textbf{snapshot}所覆盖的全部
chunk上的chunkserver lease。我们不希望这些操作运行时阻塞其他master
操作。因此，我们允许多个操作同时处于活动状态，并使用命名空间区域上的
锁来确保正确的串行化。

与许多传统文件系统不同，GFS没有一种列出目录中全部文件的、按目录划分的
数据结构\emph{(译者注：读者参考“元数据”一节的设计，它更像是维护的一个
fullpath->metadata信息的映射，并没有按目录划分的树结构)}。
它也不支持同一个文件或目录的别名(即Unix术语中的硬链接或
符号链接)。在逻辑上，GFS将命名空间表示为一张从完整路径名映射到元数据
的查找表。通过前缀压缩，这张表可以高效地存储在内存中。命名空间树中的每个节点(绝对文件名或绝对目录名)
都关联着一个读写锁。

每个master操作在运行前都会获取一组锁。通常，若操作涉及
\textbf{/d1/d2/.../dn/leaf}，它会对目录名\textbf{/d1}、\textbf{/d1/d2}、...、
\textbf{/d1/d2/.../dn}获取读锁，并对完整路径名\textbf{/d1/d2/.../dn/leaf}
获取读锁或写锁。请注意，\textbf{leaf}可以是文件，也可以是目录，具体取决于操作类型。

下面说明
这一加锁机制如何防止在将\textbf{/home/user}创建\textbf{snapshot}到\textbf{/save/user}期间，
创建文件\textbf{/home/user/foo}。
\textbf{snapshot}操作会对\textbf{/home}和\textbf{/save}获取读锁，并对\textbf{/home/user}和
\textbf{/save/user}获取写锁。文件创建操作会对\textbf{/home}和\textbf{/home/user}
获取读锁，并对\textbf{/home/user/foo}获取写锁。由于两个操作都会尝试获取
\textbf{/home/user}上相互冲突的锁，它们会被正确地串行化。
文件创建无需在父目录上获取写锁，因为不存在需要防止修改的“目录”或
类似inode的数据结构。名称上的读锁已足以防止父目录被删除。

该加锁方案的一个优点是：它允许在同一目录中并发进行mutation。例如，
可以在同一目录中并发创建多个文件：每个操作都会获取目录名上的读锁，
以及文件名上的写锁。目录名上的读锁足以防止该目录被删除、重命名或创建
\textbf{snapshot}。文件名上的写锁会将两次尝试创建同名文件的操作串行化。

由于命名空间可能
拥有许多节点，读写锁对象会按需延迟分配，并在不再使用时删除。此外，为了
防止死锁，锁会按一致的全序获取：先按命名空间树中的层级排序，同一层级内
再按字典序排序。

\subsection{replica放置}
一个GFS集群在多个层面上都具有高度分布性。它通常拥有分布在许多机器
机架上的数百个chunkserver。这些chunkserver又可能被来自同一机架或不同
机架的数百个客户端访问。位于不同机架的两台机器之间通信时，可能需要
跨越一个或多个网络交换机。
此外，一个机架的进出带宽可能低于该机架内所有机器总带宽的总和。
多层级的分布式结构给数据分布带来了独特挑战：既要保证可扩展性，又要保证可靠性和可用性。

chunk replica放置策略有两个目的：最大化数据可靠性和可用性，以及最大化网络
带宽利用率。对于这两个目标，仅将replica分散到不同机器上并不够。这样做只能
防范磁盘或机器故障，并充分利用每台机器的网络带宽。
我们还必须将chunk replica分散到不同机架上。这确保即使整个机架受损或离线，
某个chunk的部分replica仍能存活且可用。例如，机架可能因网络交换机或供电
线路等共享资源故障而离线。
这种做法还意味着，某个chunk的流量，尤其是读取流量，可以利用多个机架的
聚合带宽\emph{(译者注：前面提到了机架的带宽是较小的，因此从多个机架下载
会有更大的带宽，读者别忘了secondary机器同样可以下chunk)}。
另一方面，写入流量必须跨越多个机架传输；这是我们愿意接受的权衡。

\subsection{创建、重新复制与负载均衡}
创建chunk replica有三个原因：创建chunk、重新复制和重新均衡。

master创建chunk时，需要选择最初为空的replica应放置在哪里。它会考虑多个
因素。(1)我们希望将新replica放在磁盘空间利用率低于平均水平的
chunkserver上。长期来看，这会使各chunkserver之间的磁盘利用率趋于
均衡。
(2)我们希望限制每个chunkserver上“最近”创建的数量。虽然创建本身
开销很小，但它能可靠地预示即将到来的大量写入流量：chunk是在写入需要时
创建的；而在我们“追加一次、读取多次”的访问负载中，chunk一旦被完全
写入，通常就几乎成为只读的。
(3)如前所述，我们希望将一个chunk的replica分散到不同机架上。

一旦可用replica数量低于用户指定的目标值，master就会重新复制该chunk。
出现这种情况可能有多种原因：某个chunkserver变得不可用
\emph{(译者注：比如“lease与mutation顺序”提到的primary挂掉的场景，若心跳检测
超时，master会认为它上面的chunk不再availabel。后续如果需要一个新lease会
重新选出一个primary并且赋予lease，version随即+1，此时其实该chunk就只有两个可用replica了。
同理，如果secondary挂掉，此时master感知到了，也会标记为unavailbel，
当需要lease的时候也会进行grant new lease，通知其它chunk server：
“目前只剩下这两个replica了”。读者可以清晰地察觉到这里其实有一个lazy grant lease
的设计，即使是chunk机挂了，也不必马上重建lease，等需要的时候重建也很ok)}；
它报告自己的replica可能已损坏；其一块磁盘因错误被禁用；或者复制目标值被提高。
每个需要重新复制的chunk都会根据多个因素确定优先级。其中之一是它距离
复制目标值的差距。例如，丢失两个replica的chunk会比只丢失一个replica的chunk
获得更高优先级。此外，相比属于最近删除文件的chunk，我们优先重新复制
仍然存活的文件的chunk(section4.4)(\emph{译者注：读者可能会好奇为什么已经删除的文件
也有复制优先级，难道它也要复制？是的，因为GFS的删除其实是给一个“Deleted At”的标记然后改个名字，
到时间了再删除，这样即使文件被错误地删除了也能抢救回来})。
最后，为了尽量减少故障对正在运行的应用程序的影响，任何阻塞客户端进展
的chunk都会被提高优先级。

master选择优先级最高的chunk，并通过指示某个chunkserver直接从现有有效
replica复制chunk数据，来“克隆”它。新replica的放置目标与创建时类似：均衡
磁盘空间利用率、限制单个chunkserver上的同时进行的克隆操作，并将replica分散到
不同机架。
为了防止克隆流量挤占客户端流量，master会限制整个集群以及每个
chunkserver上的同时进行的克隆操作数量。此外，每个chunkserver都会通过限制
向源chunkserver发出的读取请求数量，来限制其为每次克隆操作使用的带宽。

最后，master会定期重新对副本负载均衡：它检查当前replica分布，并迁移replica以获得
更好的磁盘空间和负载均衡。通过这一过程，master还能逐步填充新的
chunkserver，而不是向新机器立即塞满大量新chunk及以及随之而来的大量写入流量。
新replica的放置条件与前文所述类似。此外，master还必须选择要删除的现有
replica。通常，它倾向于删除位于可用空间低于平均水平的chunkserver上的
replica，以均衡磁盘空间使用率。
\\\emph{译者注：这里很好懂，意思是一个新上任的chunkserver磁盘使用量很低
按照优先级算法，replica会成为“hotspot”，所以需要master进行渐进的、逐步的
负载均衡。}

\subsection{垃圾回收}
文件被删除后，GFS不会立即回收其占用的物理存储空间。它只会在文件和
chunk两个层级的定期垃圾回收过程中，以延迟方式回收这些空间。
我们发现，这种做法使系统更加简单，也更加可靠。

\subsubsection{机制}
当应用程序删除文件时，master会像处理其他变更一样，立即记录删除操作。
但是，它不会立即回收资源，而是仅将文件重命名为包含删除时间戳的隐藏
名称。在master定期扫描文件系统命名空间时，如果这类隐藏文件已存在超过
3天，就会将其移除(该时间间隔可配置)。在此之前，文件仍可通过新的特殊名
称读取，也可以通过将其重命名回普通名称来恢复删除。当隐藏文件
从命名空间中移除时，其内存元数据也会被清除。这实际上切断了它与
全部chunk之间的链接。

在对chunk命名空间进行类似的定期扫描时，master会识别孤立chunk(即无法
从任何文件到达的chunk)，并删除这些chunk的元数据。每个chunkserver会在
与master定期交换的HeartBeat消息中，报告自己持有的一部分chunk。
master会回复所有已不再存在于自身元数据中的chunk标识。chunkserver可以
自由删除这类chunk的replica。

\subsubsection{讨论}
在编程语言的语境中，分布式垃圾回收是一个需要复杂解决方案的难题；但在
我们的场景中，它相当简单。我们可以轻松识别全部指向chunk的引用：它们
位于完全由master维护的文件到chunk映射中。我们也能轻松识别全部chunk replica：
它们是每个chunkserver指定目录下的Linux文件。
任何master未知的这类replica都是“垃圾”。

与急切删除相比，采用垃圾回收方式回收存储空间有几个优点。第一，在组件
故障频繁的大规模分布式系统中，它简单且可靠。chunk创建可能只在部分
chunkserver上成功，从而留下master不知道存在的replica。replica删除消息
可能丢失，而master必须记得在自身或chunkserver发生故障后
重新发送这些消息。垃圾回收提供了一种统一且可靠的方式，用于清理所有
已知无用的replica。第二，它将存储回收合并到master的常规后台活动中，例如定期扫描命名空间
以及与chunkserver握手。因此，回收以批处理方式完成，成本可以被分摊。
此外，只有在master相对空闲时才执行回收，使master能够更及时地响应
需要快速处理的客户端请求。第三，延迟回收存储空间可以为意外且不可逆的删除提供一层安全保护。


根据我们的经验，主要缺点是：当存储空间紧张时，这种延迟有时会妨碍用户
精细调整存储使用。反复创建和删除临时文件的应用，可能无法立即复用已
释放的存储空间。我们通过以下方式处理这些问题：如果一个已删除文件被显式再次删除，就
加速其存储回收。我们还允许用户对命名空间的不同部分应用不同的复制和
回收策略。例如，用户可以指定某个目录树内文件的全部chunk不进行复制，并且任何被
删除的文件都立即且不可逆地从文件系统状态中移除。
\\\emph{译者注：
\\GFS的文件删除实际上就是打了一个“Deleted At”的标，并且重命名，等到三天后它才不
可逆地被master“遗忘”。在master遗忘后，就通过心跳“懒加载”地告诉chunk servers某个
文件已经删除了，chunk此时才能被chunk server从磁盘中删除。这里说的是可以选择给不同
目录设置不同的回收策略，比如某个critial文件夹5天回收；
而temp文件夹立马回收。
\\此外这里还提到了之前就已经说过的可按目录配置副本数的功能，可以给/tmp设置较少的副本数
(甚至是1)，来让它少复制甚至不复制，这样可以缓解磁盘资源紧张的问题。
}

\subsection{过期replica检测}
如果某个chunkserver发生故障，并在宕机期间遗漏了对某个chunk的
mutation，那么该chunk的replica可能变得过期(stale)。对于每个chunk，master都会
维护一个chunk版本号，用于区分最新replica和过期replica。

每当master为某个chunk grant新的lease时，都会提高该chunk的版本号，并
通知最新replica。master和这些replica都会将新版本号记录到各自的持久化状态中。
这一过程发生在通知任何客户端之前，因此也发生在客户端能够开始向该chunk
写入之前。
如果另一个replica当前不可用，它的chunk版本号就不会被推进。当该
chunkserver重启并报告其持有的chunk集合及对应版本号时，master会发现
它持有过期replica。如果master看到的版本号高于自身记录的版本号，master会
假定自己在grant lease时曾发生故障，并将较高版本视为最新版本。
\emph{
\\译者注：\\
这里特别微妙，grant lease过程中，master先在内存中生成new\_version=current\_version+1，
随即master将该“换届”信息告知给primary和secondary，这些replicas都会持久化version。
接下来master就把这个chunk version的推进持久化写入op log。请注意上述过程的任
何部位都可能宕机，因此要精细地考虑到所有边界情况。
\\在master重启后，它做的事情是：
\\1.从checkpoint+op log恢复metadata，包括namespace目录树，file->chunks mapping
以及chunk handle到chunk的mapping信息(这里信息也包括chunk version)。
\\2.从所有chunkserver上报中汇聚信息，包括chunk version和locations，补充自己的
内存结构。若遇到比自己持久化的version小的，直接视为stale；若遇到比自己version大的
直接采用；但是如果发现全部都比自己小的话，只能告诉用户这个chunk不可用了。
\\3.进入正轨，处理请求，补副本，进行垃圾回收。
\\此外，读者可以注意，此处master假如发现机器宕机，不论是primary还是secondary，都不会
立刻进行新的grant lease，它只会默默地标记这个机器unavailabel了。等下一次进行写入操作需要
lease的时候，才会进行新的grant lease——这完完全全又是一个“懒加载”的思想，避免一台机器挂掉
需要处理大量的grant lease，而导致master“惊群”。
}

master会在定期垃圾回收中删除过期replica。在此之前，当它回复客户端的
chunk信息请求时，实际上会将过期replica视为完全不存在。
作为另一层保护，master会在通知客户端哪个chunkserver持有某个chunk的
lease时，附带该chunk的版本号；或者，在克隆操作中指示某个chunkserver
从另一个chunkserver读取该chunk时，也会附带版本号。
客户端或chunkserver执行操作时会验证版本号，从而确保其访问的始终是
最新数据。